\begin{abstract}

本文从反射光的明暗条纹中估算晶圆表面薄层的厚度，并判断多次反射是否需要改变测厚结果。碳化硅与硅分别拟合，同一晶圆在不同入射角下共享厚度，避免把角度造成的光谱差异当成厚度差异。

单次返回的条纹同时受厚度和材料折射特性影响。采用一次往返场展开建立两束干涉模型，保留折射率随波数变化的色散效应，以周期投影提取条纹间距。实测条纹间距折算的周期等效量为 $19.54\,\mu\mathrm{m}$，其几何厚度换算的不确定性来自未测光学参数；合成锚例在光学、角度等条件扰动下的厚度重估范围为 $6.66$--$10.01\,\mu\mathrm{m}$。
% src:求解/问题1/结果/模型验证.json:真实附件诊断.表观光学厚度乘积_dq_微米
% src:求解/问题1/结果/灵敏度.json:条件压力范围_微米

同片碳化硅的两角谱条纹位置接近，背景与振幅却有差别，因此采用双角色散消干扰模型，让两角共享厚度和色散。算法以变投影直接求出逐角背景、振幅，再从多个起点搜索物理参数。按不贴参数边界的代表点，碳化硅厚度为 $10.80\,\mu\mathrm{m}$，连续留段、厚度与折射率组合及光学条件变化形成 $0.64$--$18.11\,\mu\mathrm{m}$ 的非概率情景范围。
% src:求解/问题2/结果/厚度结果.json:同片最佳厚度_微米;条件区间_微米;条件区间构造

多次返回是否改变测厚，以同一波段内的两束与完整往返场对照判断。采用逐次往返厚度对照，将返回光场按几何级数累加，再以多起点约束优化求解。硅的全量拟合厚度为 $2.17\,\mu\mathrm{m}$，区段、光学条件及拟合效果相近的候选厚度形成 $1.72$--$18.16\,\mu\mathrm{m}$ 的非概率范围。两材料的高阶返回收益随区段反向，碳化硅沿用前问的代表厚度与范围。
% src:求解/问题3/结果/回炉轮3候选/仲裁闭环/20260910_131501_81623_建模公式/厚度结果.json:核心指标.硅全量完整往返厚度_um;核心指标.硅已计算条件包络下限_um;核心指标.硅已计算条件包络上限_um;核心指标.碳化硅正式基准厚度_um;核心指标.碳化硅正式条件范围下限_um;核心指标.碳化硅正式条件范围上限_um;分材料验证.硅.逐折汇总.0.完整相对基线下降_%.两束;分材料验证.硅.逐折汇总.1.完整相对基线下降_%.两束;分材料验证.碳化硅.逐折汇总.0.完整相对基线下降_%.两束;分材料验证.碳化硅.逐折汇总.1.完整相对基线下降_%.两束

稳健性检验区分合成测厚误差与实测光谱预测误差：匹配合成情景中，场反演的厚度平均绝对相对误差为 $0.0008\%$，常折射率峰距法为 $2.59\%$；硅连续留段中，完整往返相对两束的标准化反射率预测误差一折降低 $25.4\%$，另一折增加 $9.52\%$。这表明谱形改善随区段变化，测厚仍依赖光学假设；缩窄厚度范围需要折射率标定，确认多光束还需核对相位关系的稳定性、光束重叠及仪器分辨能力。
% src:求解/问题1/结果/合成验证.json:主指标.数值;常折射率峰距基线.平均绝对相对误差_百分比
% src:求解/问题3/结果/回炉轮3候选/仲裁闭环/20260910_131501_81623_建模公式/厚度结果.json:分材料验证.硅.逐折汇总.0.完整相对基线下降_%.两束;分材料验证.硅.逐折汇总.1.完整相对基线下降_%.两束

\keywords{外延层厚度；双角光谱；色散；变投影；多光束干涉}
\label{abstract:end}
\end{abstract}
